\documentclass[10pt]{article}
\usepackage[margin=0.85in]{geometry}
\usepackage{amsmath,booktabs,array}
\usepackage[T1]{fontenc}
\usepackage{microtype}
\usepackage{xcolor,listings}
\usepackage[hidelinks]{hyperref}
\lstset{basicstyle=\ttfamily\footnotesize,breaklines=true,columns=fullflexible,
  keepspaces=true,showstringspaces=false,frame=single,rulecolor=\color{gray!40}}
\setlength{\parindent}{0pt}
\setlength{\parskip}{4pt}
\title{COMP 580 Assignment 2\\MinHash for Task and Content Reuse}
\author{Yifei Wang}
\date{November 2026}

\begin{document}
\maketitle
\vspace{-1.5em}

This project tests two ways to reduce repeated work in Spartacus: reuse an
answer to an equivalent task, or replace repeated tool output with a reference
and its exact differences. Combining them used 23.6\% fewer input tokens in the
sampled experiment, but the pass rate fell from 90\% to 75\%. The implementation
saved work, but the comparison did not preserve accuracy.
Code and recorded results are available at
\href{https://github.com/a-little-hoof/580-hw2}{github.com/a-little-hoof/580-hw2}.

\section{Implementation}
\textbf{MinHash and LSH.} The implementation is split into \texttt{minhash.py},
\texttt{policy.py}, and \texttt{content.py}. Text is case-folded and converted
into sets of three-word shingles, retaining numbers. Each shingle is hashed
with BLAKE2b, then mapped through 128 seeded affine functions
$h_i(x)=(a_i x+b_i)\bmod p$, where $p=2^{61}-1$. The minimum value under each
function forms one signature component. Similarity is estimated by
\[
  \widehat J(A,B)=\frac{1}{128}\sum_{i=1}^{128}
  \mathbf{1}\{s_i(A)=s_i(B)\}.
\]
The affine hashes approximate min-wise permutations. LSH groups the signature
into 32 bands of four components; a shared band retrieves a candidate for
comparison. Empty shingle sets are skipped. The hashing seed is 73.

\textbf{Task reuse.} A candidate must reach similarity $\tau_t=0.75$ and pass
checks on the fields that determine the answer. Tickets must match product,
error code, platform, and story. Bugs must match exception, top frame,
major/minor version, and the story in the Actual section. For these stories,
the policy recognizes an unambiguous cue from the public instructions;
otherwise it requires an exact match after normalization. Documents compare
identifier and status, retaining the original history ID for duplicates.
Citations compare DOI, title, authors, venue, and year. Corpus tasks compare
scenario, request kind, and target.

The cache starts with verified public history. A new answer enters it only
after passing the official verifier. Conflicting cached answers cause the
policy to decline reuse. Hidden labels are used only by the separate verifier,
never by the agent or the reuse policy.

\textbf{Content reuse.} Each task gets a fresh cache with $\tau_c=0.80$.
When a tool output resembles an earlier full output, the compressor records
line insertions, deletions, and replacements. The result identifies its source
and includes SHA-256 hashes for the base and reconstructed text. The code
checks reconstruction before returning the compressed output, preserving
numbers, identifiers, punctuation, whitespace, and line endings.

Compression is accepted only if both character count and a simple token-count
estimate fall by at least 15\%. Outputs shorter than 256 characters, errors,
and truncated outputs are returned unchanged. Compressed outputs cannot serve
as new bases, so references do not form chains. The full base stays in the
conversation; context compaction is disabled for these experiments. These
checks ensure that the original text can be recovered, but the model still
has to interpret the differences correctly.

\section{Experiments and results}
The eight datasets are tickets, bugs, docs, citations, hosts, configs, threads,
and revisions. Together they contain 420 tasks. The agent experiment uses the
first five tasks from each dataset, sorted by sequence number and task ID,
giving 40 tasks in each of four modes: baseline, task-only, content-only, and
combined. Each task/configuration was run once, for 160 completed runs.
This small sample kept API usage manageable. It is not a random sample, and
one answer changes a dataset's pass rate by 20 percentage points.

All four modes use \texttt{gemini-3.5-flash-lite}, identical prompts and tools,
temperature 0.4, and a 65,536-token output cap. Each task has a fresh workspace
containing only its public environment files. The available tools are
\texttt{read\_file}, \texttt{grep}, \texttt{list\_files}, and
\texttt{write\_file}; only \texttt{answer.json} can be written. The prompt asks
the agent to read relevant corpus files individually and explains the diff
format in every mode. The provided harness, datasets, and verifiers are
unchanged. Runs allow 20 tool-capable turns, followed by one final response
without tools if the limit is reached. No model sampling seed is set.

Accuracy is the fraction of answers accepted by the verifier. Request counts
include successful model responses, including the final response, but exclude
internal HTTP retries. Input-token counts come from the provider. Time includes
setup, verification, MinHash lookup, content processing, cache updates, and
amortized initial indexing. Requests are spaced at least five seconds apart;
the tables show this waiting time separately. Quota-interrupted attempts are
listed separately and do not count as task failures. Overnight waits between
quota resets are excluded. The earlier 41-run pilot in \texttt{results/gemini}
is also excluded; all results here come from \texttt{results/gemini-minimal}.

\begin{table}[ht]
\centering\small
\begin{tabular}{lrrrrrrr}
\toprule
Mode & Tasks & Accuracy & Tools & Requests & Input (M) & Time (h) & Pacing (h) \\
\midrule
Baseline & 40 & 90.00\% & 473 & 500 & 4.474 & 0.779 & 0.468 \\
Task-only & 40 & 85.00\% & 398 & 430 & 4.037 & 0.675 & 0.419 \\
Content-only & 40 & 87.50\% & 492 & 532 & 4.687 & 0.767 & 0.474 \\
Combined & 40 & 75.00\% & 406 & 438 & 3.419 & 0.659 & 0.354 \\
\bottomrule
\end{tabular}
\caption{Sampled evaluation: the first 5 incoming tasks per dataset, sorted by (seq, tid), with the same selection in all four modes. Time includes setup, verification, reuse processing, and API pacing. Requests exclude HTTP retries.}
\label{tab:online-aggregate}
\end{table}
\begin{table}[ht]
\centering\small
\begin{tabular}{lrrrr}
\toprule
Dataset & Baseline & Task-only & Content-only & Combined \\
\midrule
tickets & 2/5 (40.0\%) & 2/5 (40.0\%) & 3/5 (60.0\%) & 1/5 (20.0\%) \\
bugs & 5/5 (100.0\%) & 4/5 (80.0\%) & 5/5 (100.0\%) & 4/5 (80.0\%) \\
docs & 5/5 (100.0\%) & 5/5 (100.0\%) & 5/5 (100.0\%) & 4/5 (80.0\%) \\
citations & 5/5 (100.0\%) & 3/5 (60.0\%) & 2/5 (40.0\%) & 3/5 (60.0\%) \\
hosts & 5/5 (100.0\%) & 5/5 (100.0\%) & 5/5 (100.0\%) & 5/5 (100.0\%) \\
configs & 5/5 (100.0\%) & 5/5 (100.0\%) & 5/5 (100.0\%) & 4/5 (80.0\%) \\
threads & 5/5 (100.0\%) & 5/5 (100.0\%) & 5/5 (100.0\%) & 5/5 (100.0\%) \\
revisions & 4/5 (80.0\%) & 5/5 (100.0\%) & 5/5 (100.0\%) & 4/5 (80.0\%) \\
\bottomrule
\end{tabular}
\caption{Correct answers according to the supplied verifier. Interrupted attempts are excluded.}
\label{tab:online-accuracy}
\end{table}
\begin{table}[ht]
\centering\small
\begin{tabular}{lrrrrr}
\toprule
Dataset & Baseline & Task-only & Content-only & Combined & Saved \\
\midrule
tickets & 491,641 & 642,585 & 800,850 & 562,401 & -14.39\% \\
bugs & 281,620 & 484,597 & 461,424 & 379,609 & -34.79\% \\
docs & 393,175 & 248,573 & 655,901 & 192,207 & 51.11\% \\
citations & 1,247,466 & 913,176 & 924,426 & 795,951 & 36.19\% \\
hosts & 404,858 & 478,838 & 295,284 & 339,115 & 16.24\% \\
configs & 642,756 & 138,507 & 491,648 & 356,558 & 44.53\% \\
threads & 354,643 & 598,721 & 513,454 & 341,010 & 3.84\% \\
revisions & 657,431 & 531,919 & 544,478 & 451,757 & 31.28\% \\
\bottomrule
\end{tabular}
\caption{Input tokens reported by Gemini. Saved compares combined reuse with baseline; negative values mean more tokens were used.}
\label{tab:online-tokens}
\end{table}
\begin{table}[ht]
\centering\small
\begin{tabular}{lrrrrr}
\toprule
Mode & Interruptions & Requests & Input tokens & Time (h) & Pacing (h) \\
\midrule
Baseline & 3 & 20 & 147,348 & 0.058 & 0.020 \\
Task-only & 0 & 0 & 0 & 0.000 & 0.000 \\
Content-only & 4 & 41 & 577,708 & 0.095 & 0.037 \\
Combined & 0 & 0 & 0 & 0.000 & 0.000 \\
\bottomrule
\end{tabular}
\caption{Extra work from interrupted attempts, reported separately from completed tasks. Time includes pacing and retries, but not the wait between restarts.}
\label{tab:online-interruptions}
\end{table}
Including interrupted attempts, the run used 1,961 model responses and 17,341,635 input tokens over 3.033 hours (1.771 hours of API pacing).


\textbf{Agent results.} Task-only reuse reduced model requests from 500 to 430
and input tokens by 9.8\%, while passing 34/40 tasks versus the baseline's
36/40. Content-only reuse used 4.8\% more tokens overall and passed 35/40.
Combined reuse used the fewest tokens, 3.419 million, but passed only 30/40.
Its lower cost therefore came with a lower observed pass rate.

The added processing time was small: 1.214 seconds for task-only, 3.408 seconds
for content-only, and 4.001 seconds for combined reuse across their 40 tasks.
These totals include lookup, cache updates, initial indexing, and content
processing. Most recorded wall time came from model calls and request pacing.
With only one run per configuration, this experiment cannot separate the effect
of reuse from variation in model responses.

\textbf{What went wrong?} Comparing the same tasks, eight baseline passes became
failures in combined mode, while two baseline failures became passes. This
accounts for the net loss of six correct answers. None of the eight regressions
used a cached answer. Five also had no compressed tool output: their traces
show missed task cues or overlooked differences from similar records.
For example, \texttt{T-52023e} supplied the author name ``Qiunn,'' while the
database said ``Quinn.'' The combined run read both spellings but reported
only the incorrect year. That failure cannot be explained by text removed by
the compressor.

\begin{table}[ht]
\centering\small
\begin{tabular}{p{0.26\linewidth}rp{0.60\linewidth}}
\toprule
Task & Diffs & Observed failure \\
\midrule
tickets \texttt{T-a1baca} & 0 & Chose a backup fix for a network-drive issue. \\
tickets \texttt{T-53f313} & 0 & Chose an old-release fix for a clock issue. \\
bugs \texttt{T-8d0531} & 0 & Chose a deadlock cause instead of the upgrade cause. \\
docs \texttt{T-bec0fe} & 4 & Repeated searches; reached the turn limit without writing an answer. \\
citations \texttt{T-3f588b} & 0 & Missed ``recognition'' versus ``detection'' in the title. \\
citations \texttt{T-52023e} & 0 & Found the wrong year but missed a misspelled author. \\
configs \texttt{T-b06652} & 6 & Invalid searches; reached the turn limit without writing an answer. \\
revisions \texttt{T-e44ab3} & 4 & Returned \texttt{v07}, although R9 was present in \texttt{v05}. \\
\bottomrule
\end{tabular}
\caption{The eight tasks that passed in baseline and failed in combined mode.
Diffs counts compressed tool outputs in the combined trace.}
\label{tab:failures}
\end{table}

The three runs that did receive compressed outputs failed in different ways.
Docs and configs exhausted the 20 tool-capable turns without calling
\texttt{write\_file}; their null answer artifacts mean no answer was produced.
Repeated or invalid searches were visible before the first compression in
both runs. The config trace even contained the violating value
\texttt{auth.allow\_anonymous: true} in a diff, but never submitted it.
The revisions run wrote \texttt{v07/R9} after receiving a diff for \texttt{v05}
that retained R9 from its full base. Comparing \texttt{v04} and \texttt{v05}
places the addition at \texttt{v05}.

All 14 diffs in these three traces reconstruct exactly. The audit also
replays their recorded tool calls on copies of the public files and checks
the reconstructed outputs. This rules out missing text in these diffs, but
does not rule out difficulty interpreting the reference format. The evidence
supports separating wrong answer selection, unsuccessful tool use, and
compression handling rather than treating the whole accuracy drop as one
failure of MinHash. A controlled rerun with full text would be needed to test
whether the format contributed to the last three failures.

\textbf{Offline checks.} Two larger checks help separate cache behavior from
model behavior. The task check queries the original verified history for all
420 tasks and verifies every proposed reuse. The history stays fixed throughout
the check. At $\tau_t=0.75$, 79 tasks (18.81\%) reuse an answer, and all 79 pass.
Across the six thresholds in Table~\ref{tab:threshold}, all 462 accepted reuse
decisions pass; some tasks appear at several thresholds.

The content check reads each corpus scenario's files in sorted order, using
Spartacus's line numbering and truncation rules. Each scenario starts with an
empty cache. All 349 compressed outputs reconstruct exactly. These savings
measure tool-output characters, whereas the agent experiment measures actual
model input tokens.

\begin{table}[ht]
\centering\small
\begin{tabular}{lrrrrr}
\toprule
Dataset & Tasks & Reused & Passed & Build (ms) & Lookup (ms/task) \\
\midrule
tickets & 70 & 6 & 6 & 221.7 & 1.47 \\
bugs & 60 & 10 & 10 & 258.3 & 1.92 \\
docs & 60 & 10 & 10 & 261.8 & 2.45 \\
citations & 70 & 24 & 24 & 69.4 & 0.42 \\
hosts & 40 & 10 & 10 & 65.6 & 1.00 \\
configs & 40 & 6 & 6 & 71.6 & 0.99 \\
threads & 40 & 6 & 6 & 76.4 & 0.99 \\
revisions & 40 & 7 & 7 & 72.8 & 1.07 \\
\midrule
Total & 420 & 79 & 79 & -- & -- \\
\bottomrule
\end{tabular}
\caption{History-only offline reuse at $\tau_t=0.75$. Passed counts only accepted reuse decisions, not all incoming tasks. Build and lookup timings exclude verifier execution.}
\label{tab:tasks}
\end{table}
\begin{table}[ht]
\centering\small
\begin{tabular}{rrrrrrr}
\toprule
$\tau_t$ & 0.50 & 0.65 & 0.75 & 0.85 & 0.95 & 1.00 \\
\midrule
Reused & 96 & 90 & 79 & 71 & 63 & 63 \\
Verifier failures & 0 & 0 & 0 & 0 & 0 & 0 \\
\bottomrule
\end{tabular}
\caption{Offline threshold sweep over the same 420 incoming tasks and fixed histories; no fresh answers are indexed. The sweep measures reuse coverage, not full-agent accuracy.}
\label{tab:threshold}
\end{table}
\begin{table}[ht]
\centering\small
\begin{tabular}{lrrrrr}
\toprule
Dataset & Outputs & Deltas & Original chars & Saved (\%) & Process (s) \\
\midrule
hosts & 175 & 83 & 439,877 & 22.97 & 1.59 \\
configs & 212 & 142 & 459,692 & 38.31 & 1.19 \\
threads & 710 & 0 & 681,171 & 0.00 & 2.27 \\
revisions & 198 & 124 & 434,247 & 34.29 & 1.38 \\
\bottomrule
\end{tabular}
\caption{Offline corpus replay at $\tau_c=0.80$, 24 scenarios per dataset, each with a fresh cache. All 349 emitted deltas reconstruct exactly. Savings are tool-output characters, not measured model tokens.}
\label{tab:content}
\end{table}


\section{Discussion}
\textbf{1. When is task reuse beneficial?}
It helps most when the same question returns and a verified answer is already
available. In the agent experiment, task-only reuse skipped eight tasks and
combined reuse skipped seven; every reused answer passed, with no model
requests or tool calls for those tasks. In the larger history check, citations
had 24/70 reuse hits, compared with 6/70 for tickets. Citation fields give the
policy direct checks for equivalence; ticket stories require more interpretation.

\textbf{2. Is high similarity sufficient for safe reuse?}
No. Bug-history records \texttt{H-84318e} and \texttt{H-fff543} both describe
a \texttt{PermissionError} at \nolinkurl{lattice.store.blob.write} caused by
missing directory ownership. Their signatures agree in 119 of 128 positions,
giving similarity 0.930. However, their versions are 3.2.1 and 3.3.1, and their
verified answers are \texttt{storage/RC-05} and \texttt{storage/RC-14}.
The version check rejects the match despite its high similarity.

\textbf{3. How does the threshold affect performance?}
Raising the threshold reduced offline reuse from 96/420 tasks at 0.50 to
63/420 at 1.00. None of the accepted answers failed verification in this sweep.
Here, the lower threshold found more usable matches because the field checks
still rejected unsafe candidates. The default 0.75 was set before the sweep.
A single threshold is unlikely to suit every dataset: a citation and a long
bug report have very different wording. LSH can also miss a candidate before
the threshold is applied.

\textbf{4. When is content-level reuse effective?}
It works well when long files share most of their text and differ in a few
important lines. On hosts, content-only reuse reduced input tokens from
404,858 to 295,284 (27.1\%), with 5/5 passes in both modes. The offline replay
also found substantial character savings for configs (38.31\%) and revisions
(34.29\%). Threads behaved differently: the offline replay compressed none of
its 710 outputs, and the content-only agent used 44.8\% more tokens than the
baseline. Short or heavily changed outputs often fail the savings check, and
different agent reads or extra turns can outweigh savings from individual diffs.

\textbf{5. Which optimization is more valuable?}
Task reuse is the better starting point when verified duplicates are common:
one hit avoids the entire agent run. Content reuse is useful when a new task
still needs to read similar files. In this sample, task-only reuse reduced both
requests and tokens, while content-only reuse increased both overall. Combined
reuse saved the most tokens but had the lowest accuracy, so these results do
not support using it without further checking.

The modes can also build different caches. For example, task-only solved config
task \texttt{T-b06652} and reused its verified answer for \texttt{T-025f77}.
Combined reuse failed the first task, so that answer could not enter its cache
and the second task needed another agent run. The trace audit shows that the
first run exhausted its turns without writing an answer. This explains one
difference in reuse counts without implying an incorrect cached answer.

\textbf{6. What are MinHash's limitations?}
MinHash measures shared wording, not whether two tasks have the same answer.
Paraphrases can look unrelated, while a one-number change can leave similarity
almost unchanged. The bug-version example shows why extra checks are needed.
Three-word sets also lose order beyond each shingle, and the field checks here
depend on the assignment's schemas. For short outputs, computing a signature
and describing a diff may cost more than sending the original text. One useful
improvement would be to split the index by critical fields, such as product and
version, before running LSH. That would reduce irrelevant candidates while
keeping the final validation checks.

\section{Reproduction}
The code requires Python 3.9 or newer and Git, with no third-party Python
packages. The assignment and Spartacus submodules are pinned to
\texttt{a7c5076} and \texttt{711db62}, respectively; the full commits are recorded
in the run settings. To repeat the checks in new output directories:
\begin{lstlisting}
git clone --recurse-submodules git@github.com:a-little-hoof/580-hw2.git
cd 580-hw2
python -m unittest discover -s tests -v
python evaluate.py --offline --output results/reproduction-offline
\end{lstlisting}
For the model experiment, set \texttt{GEMINI\_API\_KEY} in the environment or
in \nolinkurl{~/.config/spartacus/env}. Keep the key outside the repository.
\begin{lstlisting}
python evaluate.py --output results/reproduction-gemini \
  --datasets tickets bugs docs citations hosts configs threads revisions \
  --modes baseline task content combined \
  --model gemini-3.5-flash-lite --max-turns 20 --limit 5 \
  --task-threshold 0.75 --content-threshold 0.80 \
  --request-interval 5
\end{lstlisting}
If interrupted, repeat the same command with \texttt{--resume} appended.
Resume requires unchanged settings and source files and restores the verified
cache from completed tasks. To rebuild the tables from the submitted results:
\begin{lstlisting}
python report/summarize.py results/gemini-minimal
python report/analyze_failures.py results/gemini-minimal
\end{lstlisting}
The summary script checks that every selected task has one result in each
mode before writing the tables. The failure audit writes
\nolinkurl{report/failure_analysis.json}, including paired outcomes and
one-based transcript message indices. It uses saved traces and public files;
no model requests or hidden labels are needed. Raw answers, transcripts, and
interruption records are included alongside the metrics.
\end{document}
